%! Introducing Statistics: What Can We Learn from Data? — Unit 1, Lesson 1.0 — Warm-Up (Answer Key)
\documentclass[10pt]{article}
\usepackage{apstats-article}
\usepackage{apstats-key}   % pulls in -boxes; do NOT also load -boxes

\begin{document}

\pageheader{Unit 1, Lesson 1.0}{Warm-Up --- Answer Key}

\vspace{0.15in}

\boxguard[24]
\begin{hookbox}
\small
Here are seven numbers. That is all you are given.

\vspace{8pt}
\begin{center}
{\large\bfseries\color{navy} 68 \qquad 71 \qquad 64 \qquad 70 \qquad 66 \qquad 73 \qquad 69}
\end{center}
\vspace{4pt}
\end{hookbox}

\vspace{0.16in}

\begin{enumerate}[leftmargin=1.7em, itemsep=0.16in]

  \item What do these numbers tell you? Answer honestly.\\[4pt]
        \ansline{Nothing yet. They could be test scores, temperatures, ages, or heights ---}\\[1pt]\ansline{a number with no context carries no information at all.}

  \item List \textbf{three} things you would need to know before these numbers could
        mean anything.\\[4pt]
        \ansline{1. WHO was measured --- which individuals these seven numbers describe.}\\[1pt]\ansline{2. WHAT was measured --- the variable, such as height, score, or temperature.}\\[1pt]\ansline{3. The UNITS --- inches, degrees, or points --- and when and where it was taken.}

  \item One of these questions can be answered with a single measurement. The other
        cannot. \textbf{Circle} the one that needs many measurements, then explain why.
        \vspace{6pt}

        \hspace{1.2em}\textbf{(A)} How tall is Maria?
        \hfill
        \textbf{(B)} How tall are the students in this class?
        \textcolor{keyred}{\textbf{$\leftarrow$ circle}}
        \vspace{8pt}

        Why?\\[4pt]
        \ansline{(B). Maria has exactly one height, so one measurement answers (A) fully.}\\[1pt]\ansline{Students' heights differ from person to person, so (B) needs many values.}

\end{enumerate}

\end{document}
